% Overleaf compiler: pdfLaTeX
\documentclass[11pt,a4paper]{article}
\usepackage[utf8]{inputenc}
\usepackage[T1]{fontenc}
\usepackage[margin=2.5cm]{geometry}
\usepackage{amsmath,amssymb}
\usepackage{graphicx}
\usepackage{booktabs}
\usepackage{hyperref}

\title{Notes on Heat Flow in a Thin Rod}
\author{A. Writer}
\date{\today}

\begin{document}
\maketitle

\begin{abstract}
We derive the one-dimensional heat equation, solve it by separation of variables
and compare the first few modes.
\end{abstract}

\section{Introduction}\label{sec:intro}
Heat flows from hot to cold. In a thin rod of length $L$ the temperature $u(x,t)$
depends on one space variable only.\footnote{We ignore losses through the surface.}
Section~\ref{sec:model} sets up the model and Section~\ref{sec:solution} solves it.

\section{The model}\label{sec:model}
Fourier's law and conservation of energy give
\begin{equation}\label{eq:heat}
  \frac{\partial u}{\partial t} = \alpha \frac{\partial^2 u}{\partial x^2}, \qquad 0 < x < L,\ t > 0,
\end{equation}
with $u(0,t) = u(L,t) = 0$ and $u(x,0) = f(x)$.

\section{Solution}\label{sec:solution}
Trying $u(x,t) = X(x)T(t)$ in \eqref{eq:heat} gives
\begin{align*}
  X'' + \lambda X &= 0, \\
  T' + \alpha\lambda T &= 0,
\end{align*}
so that
\[
  u(x,t) = \sum_{n=1}^{\infty} b_n \sin\left(\frac{n\pi x}{L}\right) e^{-\alpha (n\pi/L)^2 t}.
\]

The steps are:
\begin{enumerate}
  \item separate the variables;
  \item solve the eigenvalue problem for $X$;
  \item match the initial condition with a Fourier sine series.
\end{enumerate}

\begin{table}[h]
  \centering
  \begin{tabular}{lcc}
    \toprule
    Mode $n$ & Decay rate & Half-life \\
    \midrule
    1 & $\alpha\pi^2/L^2$  & $0.070\,L^2/\alpha$ \\
    2 & $4\alpha\pi^2/L^2$ & $0.018\,L^2/\alpha$ \\
    \bottomrule
  \end{tabular}
  \caption{The first two modes.}\label{tab:modes}
\end{table}

\section{Conclusion}
Higher modes die out fast, as Table~\ref{tab:modes} shows. See
\url{https://en.wikipedia.org/wiki/Heat_equation} for more.

\end{document}
